\documentclass[11pt,a4paper]{article}

% ---------- Pacotes ----------
\usepackage[utf8]{inputenc}
\usepackage[T1]{fontenc}
\usepackage[portuguese]{babel}
\usepackage[a4paper,margin=2.5cm]{geometry}
\usepackage{titlesec}
\usepackage{xcolor}
\usepackage{enumitem}
\usepackage{tcolorbox}
\usepackage{hyperref}
\usepackage{fancyhdr}
\usepackage{graphicx}

\tcbuselibrary{skins,breakable}

% ---------- Cores ----------
\definecolor{corprincipal}{RGB}{31,73,125}
\definecolor{corsecundaria}{RGB}{68,114,148}
\definecolor{cinzaclaro}{RGB}{245,245,245}

% ---------- Formatação de secções ----------
\titleformat{\section}
  {\normalfont\Large\bfseries\color{corprincipal}}
  {\thesection}{0.8em}{}
\titleformat{\subsection}
  {\normalfont\large\bfseries\color{corsecundaria}}
  {\thesubsection}{0.7em}{}

% ---------- Cabeçalho / rodapé ----------
\pagestyle{fancy}
\fancyhf{}
\renewcommand{\headrulewidth}{0.4pt}
\fancyhead[L]{\small Sprint Timer -- Relatório de Projeto}
\fancyhead[R]{\small José Oliveira}
\fancyfoot[C]{\thepage}

% ---------- Caixa de problema/solução ----------
\newtcolorbox{probsol}[1]{
  breakable, enhanced,
  colback=cinzaclaro, colframe=corsecundaria,
  boxrule=0.5pt, arc=2pt,
  left=6pt, right=6pt, top=4pt, bottom=4pt,
  title={#1}, fonttitle=\bfseries\color{white},
  colbacktitle=corsecundaria,
}

% ---------- Placeholder de imagem ----------
\newcommand{\imagemplaceholder}[2]{
  \begin{center}
  \fbox{
    \begin{minipage}[c][#1][c]{0.85\textwidth}
      \centering
      \color{gray}\itshape #2
    \end{minipage}
  }
  \end{center}
}

\hypersetup{
    colorlinks=true,
    linkcolor=corprincipal,
    urlcolor=corprincipal,
}

\title{}
\author{}
\date{}

\begin{document}

% ================= PÁGINA DE TÍTULO =================
\begin{titlepage}
\centering
\vspace*{3cm}

{\Huge\bfseries\color{corprincipal} Raul 3000 - Sprint Timer}\\[0.5cm]
{\Large Medidor de Velocidade de Sprint com Sensores de Movimento}\\[2cm]

{\large Relatório de Projeto}\\[3cm]

\rule{0.6\textwidth}{0.4pt}\\[0.5cm]
{\large José Oliveira}\\[0.3cm]
{\normalsize \today}

\vfill

{\small Projeto de eletrónica e automação}

\end{titlepage}

% ================= INTRODUÇÃO =================
\section*{Introdução}
\addcontentsline{toc}{section}{Introdução}

O projeto \emph{Raul 3000 - Sprint Timer} tem como objetivo medir a velocidade de um sprint através do tempo decorrido entre a ativação de dois sensores de movimento: um sensor de partida e um sensor de chegada. O tempo medido é apresentado ao utilizador num ecrã LCD integrado numa caixa de plástico, permitindo uma utilização simples e portátil.

O sistema é construído em torno de um microcontrolador Raspberry Pi Pico, responsável por gerir os sensores, contar o tempo e controlar o ecrã de resultados. Ao longo do desenvolvimento, o projeto passou por duas iterações principais: um primeiro protótipo funcional, alimentado por um pack de baterias não recarregável, e um protótipo final, que introduz uma solução de carregamento integrada e monitorização do estado da bateria.

Os principais componentes utilizados no projeto são:

\begin{itemize}[itemsep=1pt]
    \item Sensores de movimento (2)
    \item Relays mecânicos (2)
    \item Microcontrolador Raspberry Pi Pico (1)
    \item Módulo DC-DC \textit{step-up} (1)
    \item Módulo DC-DC \textit{step-down} (1)
    \item Módulo de carregamento com BMS (1)
    \item Display LCD (1)
    \item Módulo adaptador I2C (1)
    \item Baterias de lítio 3.7V (2)
    \item Caixa de plástico (1)
    \item Cabos diversos
\end{itemize}

Este relatório descreve o percurso de desenvolvimento do projeto, desde o conceito inicial até ao protótipo final, apresentando os problemas encontrados em cada fase e as soluções adotadas para os resolver.

% ================= SECÇÃO 1: CONCEITO INICIAL =================
\section{Conceito Inicial}

A ideia inicial do projeto assentava num sistema simples, composto por dois sensores de movimento e um ``cérebro'' central. Um dos sensores funcionaria como sensor de partida, iniciando a contagem do tempo, enquanto o segundo funcionaria como sensor de chegada, terminando essa contagem. A diferença de tempo entre a ativação dos dois sensores corresponderia, na prática, ao tempo de sprint entre eles.

Este conceito serviu de base para todo o desenvolvimento posterior, sendo que a principal evolução do projeto consistiu em definir como implementar fisicamente este ``cérebro'' e como resolver os desafios de compatibilidade elétrica e de alimentação que foram surgindo.

% ================= SECÇÃO 2: PRIMEIRO PROTÓTIPO =================
\section{Primeiro Protótipo Funcional}

Nesta fase do desenvolvimento, o objetivo foi construir um protótipo totalmente funcional, ainda sem uma solução de alimentação integrada e recarregável. Foram identificados e resolvidos cinco problemas principais.

\subsection{Protótipo 0}

Antes de se chegar a uma versão funcional dentro de caixa, o circuito foi montado e testado numa placa de desenvolvimento (\textit{breadboard}), permitindo validar as ligações entre o microcontrolador, os sensores e os restantes componentes antes da integração final. Por facilidade de referência ao longo deste relatório, esta fase de testes iniciais é designada como ``Protótipo 0''.

\imagemplaceholder{6cm}{Fotografia do circuito montado em placa de desenvolvimento (Protótipo 0).}

\subsection{O Cérebro do Sistema}

\begin{probsol}{Problema}
Como implementar o ``cérebro'' do sistema e como o ligar aos dois sensores de movimento.
\end{probsol}

\begin{probsol}{Solução}
Foi utilizado um microcontrolador Raspberry Pi Pico, programado com um código simples em Python. O Pico está diretamente ligado aos sensores por cabo: inicia a contagem do tempo quando o sensor de partida é ativado e termina a contagem quando o sensor de chegada é ativado.
\end{probsol}

\subsection{Compatibilização de Tensões}

\begin{probsol}{Problema}
Os sensores de movimento operam entre 12V e 24V, enquanto o Raspberry Pi Pico apenas suporta tensões entre 3.3V e 5V. Não é possível ligá-los diretamente.
\end{probsol}

\begin{probsol}{Solução}
Foram utilizados relays mecânicos para separar fisicamente as correntes dos sensores e do microcontrolador, permitindo ainda assim a comunicação entre ambos. Esta solução implica a existência de duas fontes de alimentação com voltagens diferentes.
\end{probsol}

\subsection{Visualização dos Resultados}

\begin{probsol}{Problema}
Como permitir ao utilizador visualizar o resultado das corridas medidas.
\end{probsol}

\begin{probsol}{Solução}
Foi integrado um pequeno ecrã LCD, que apresenta os três últimos tempos registados (o mais recente e os dois anteriores, evitando a perda de informação em caso de uma leitura falsa) e o estado atual do sistema.
\end{probsol}

O sistema opera segundo quatro estados distintos:

\begin{itemize}
    \item \textbf{A iniciar} -- Ao ligar o projeto, este demora 5 segundos a arrancar e a passar para o estado seguinte. Durante este período é possível verificar que o ecrã está a funcionar corretamente, mas não é possível iniciar uma contagem, mesmo que o sensor de partida seja ativado.
    \item \textbf{Pronto para medição} -- Após o estado inicial, ou após a apresentação de uma medição, o sistema fica pronto para iniciar um novo sprint. Basta ativar o sensor de partida para começar a medição.
    \item \textbf{A medir / em funcionamento} -- Após a ativação do sensor de partida (estando o sistema no estado ``pronto para medição''), inicia-se a contagem do tempo, que só termina quando o sensor de chegada é ativado. Este estado só é alcançado se o estado anterior for ``pronto para medição''; a ativação do sensor de partida em qualquer outro estado, incluindo este, não inicia uma nova medição.
    \item \textbf{A mostrar medição / resultado} -- Alcançado quando o sensor de chegada é ativado durante uma medição em curso. Neste estado, o sistema apresenta o último resultado durante 5 segundos, não sendo possível iniciar uma nova medição nesse intervalo.
\end{itemize}

\imagemplaceholder{6cm}{Imagem da topologia do projeto: Raspberry Pi Pico, os dois sensores ligados através dos relays, o módulo step-down, o ecrã LCD e a fonte de alimentação.}

\imagemplaceholder{5cm}{Fotografia do ecrã LCD (ainda sem tempos registados).}

\subsection{Alimentação do Sistema}

\begin{probsol}{Problema}
Como alimentar o projeto de forma prática e portátil.
\end{probsol}

\begin{probsol}{Solução}
Foram consideradas duas opções: alimentar o projeto através de cabo, ligado a uma fonte de alimentação de corrente alternada (tomada), ou dotar o projeto de uma fonte de alimentação própria, integrada. Foi escolhida a segunda opção, através de um pack de baterias de lítio integrado.
\end{probsol}

\subsection{Fontes de Alimentação Distintas a Partir de um Único Pack}

\begin{probsol}{Problema}
Como obter duas fontes de alimentação com voltagens diferentes a partir do mesmo pack de baterias, uma vez que os sensores e o microcontrolador requerem tensões distintas.
\end{probsol}

\begin{probsol}{Solução}
A solução inicial consistiu num pack de 6 baterias de 3.7V ligadas em série, resultando numa tensão de aproximadamente 22.2V. O pack de baterias foi ligado diretamente aos sensores, enquanto para o Raspberry Pi Pico foi utilizado um módulo DC-DC \textit{step-down}, configurado para fornecer entre 4V e 5V.
\end{probsol}

Com esta solução, o projeto ficou totalmente funcional. No entanto, permanecia um problema por resolver: o pack de baterias tinha de ser periodicamente substituído ou recarregado externamente, não existindo nesta iteração uma solução de carregamento integrada.

% ================= SECÇÃO 3: PROTÓTIPO FINAL =================
\section{Protótipo Final}

O segundo e último protótipo tem como principal objetivo resolver o problema de alimentação deixado em aberto na iteração anterior, introduzindo uma solução de carregamento integrada.

\subsection{Solução de Carregamento Integrada}

\begin{probsol}{Problema}
O pack de baterias do primeiro protótipo não podia ser recarregado de forma integrada, exigindo substituição ou carregamento externo periódico.
\end{probsol}

\begin{probsol}{Solução}
Foi introduzido um módulo de carregamento com proteção de baterias integrado (BMS). Para viabilizar esta solução, foi necessário reduzir o pack de baterias de 6 para 2 unidades, resultando numa tensão de aproximadamente 7.4V. Como esta tensão deixou de ser suficiente para alimentar diretamente os sensores, foi introduzido um módulo DC-DC \textit{step-up}, elevando a tensão para aproximadamente 20V.
\end{probsol}

\imagemplaceholder{6cm}{Imagem da topologia do protótipo final, incluindo o módulo de carregamento BMS e o módulo step-up.}

\imagemplaceholder{5cm}{Fotografia da caixa a carregar.}

\subsection{Monitorização do Estado da Bateria}

Com a introdução de um pack de baterias recarregável, tornou-se necessário permitir ao utilizador monitorizar o estado de carga da bateria. Para isso, foi adicionada ao ecrã uma amostragem da percentagem de bateria disponível.

\begin{probsol}{Problema}
O Raspberry Pi Pico não está ligado diretamente à bateria, mas sim a um módulo DC-DC \textit{step-down}, pelo que recebe sempre uma tensão constante de aproximadamente 5V, independentemente do estado real da bateria (ou 0V caso esta esteja descarregada). Isto impossibilita a leitura direta do nível de carga.
\end{probsol}

\begin{probsol}{Solução}
Foi utilizado um divisor de tensão para ligar o Raspberry Pi Pico diretamente à bateria através de um pino de conversão analógico-digital, permitindo a leitura da tensão real da bateria e, consequentemente, o cálculo da sua percentagem de carga.
\end{probsol}

\imagemplaceholder{5cm}{Fotografia do circuito do divisor de tensão.}

\subsection{Integração e Montagem Final}

Com todos os componentes e funcionalidades resolvidos, o circuito foi integrado dentro da caixa de plástico, resultando no produto final do projeto.

\imagemplaceholder{6cm}{Fotografia do interior da caixa, já com todos os componentes montados.}

\imagemplaceholder{6cm}{Fotografia da caixa vista de frente.}

\imagemplaceholder{6cm}{Fotografia do projeto em funcionamento.}

% ================= CONCLUSÃO =================
\section*{Conclusão}
\addcontentsline{toc}{section}{Conclusão}

O projeto \emph{Sprint Timer} evoluiu de um conceito simples, baseado em dois sensores de movimento e um microcontrolador, para um sistema portátil e autónomo capaz de medir e apresentar a velocidade de um sprint com fiabilidade. Ao longo do desenvolvimento foram resolvidos desafios de compatibilidade elétrica entre sensores e microcontrolador, de gestão de estados do sistema e, por fim, de alimentação energética.

A evolução do primeiro para o protótipo final traduziu-se, sobretudo, na passagem de uma fonte de alimentação não recarregável para uma solução de carregamento integrada com monitorização do estado da bateria, tornando o sistema mais prático e autónomo para uma utilização continuada.

% ================= ANEXO =================
\newpage
\section*{Anexo -- Fotografias Adicionais do Desenvolvimento}
\addcontentsline{toc}{section}{Anexo -- Fotografias Adicionais do Desenvolvimento}

Esta secção reúne fotografias adicionais do processo de desenvolvimento do Protótipo 0, documentadas ao longo dos testes iniciais do circuito.

\imagemplaceholder{5cm}{Fotografia do Protótipo 0 em placa de desenvolvimento (vista de cima).}

\imagemplaceholder{5cm}{Fotografia do Protótipo 0 em placa de desenvolvimento.}

\imagemplaceholder{5cm}{Fotografia do Protótipo 0 junto ao caderno de anotações do desenvolvimento.}

\end{document}